\documentclass[10pt,letterpaper]{article}
\usepackage[utf8]{inputenc}
\usepackage[spanish,es-noshorthands]{babel}
\usepackage{amsmath,amssymb}
\usepackage{graphicx}
\usepackage[american]{circuitikz}
\usepackage[letterpaper, top=1in, bottom=1in, left=1in, right=1in]{geometry}

% componentes algo mas compactos que el valor por defecto
\ctikzset{bipoles/length=1.15cm}
% signo + junto a la placa larga (borne positivo) de cada fuente
\newcommand{\bplus}[1]{\node[font=\footnotesize, inner sep=0pt] at (#1) {$+$};}
% respuesta: "R: valor [unidad]" sin cortes de linea
\newcommand{\resp}[1]{\mbox{\textbf{R:}~#1}}

\renewcommand{\labelenumi}{\arabic{enumi}.}
\setlength{\parindent}{0pt}
\pagestyle{plain}

\begin{document}

\begin{center}
    \includegraphics[height=2.2cm]{logo ubb.png}\\[6pt]
    \textbf{Electromagnetismo}\\[2pt]
    \emph{Gu\'ia N$^{\circ}$8 - Circuitos el\'ectricos en corriente continua}\\[2pt]
    \textit{Profesor: Arturo Fern\'andez P\'erez}
\end{center}

\medskip
En todos los circuitos las fuentes de fem son ideales (sin resistencia interna) y el signo $+$
se\~nala su borne positivo (placa larga del s\'imbolo). Los valores de las resistencias y de las
fem se indican en cada figura.

\begin{enumerate}
\setlength{\itemsep}{12pt}

%==================================================================
\item \begin{minipage}[t]{\linewidth}
Para el circuito de la Fig.~1, use las leyes de Kirchhoff para determinar la corriente $I$ que
circula por $R_5$, con el sentido de referencia indicado por la flecha (un resultado negativo
indica que la corriente circula en sentido contrario). \resp{$I = 2$~[A]}
\begin{center}
\begin{circuitikz}[every node/.style={font=\small}]
  % rama izquierda: eps1 (+ arriba) y R5 (flecha I hacia arriba)
  \draw (0,3.5) to[battery1, l_={$\varepsilon_1 = 20\,\text{V}$}] (0,1.9);
  \draw (0,0) to[R, l={$R_5 = 3\,\Omega$}, f>_={$I$}] (0,1.9);
  \bplus{0.32,2.95}
  % rama superior
  \draw (0,3.5) to[R, l=$R_1$, a=$4\,\Omega$] (3.5,3.5)
                to[R, l=$R_2$, a=$6\,\Omega$] (7,3.5);
  % rama central y derecha
  \draw (3.5,0) to[R, l=$R_3$, a=$12\,\Omega$] (3.5,3.5);
  \draw (7,0) to[R, l=$R_4$, a=$2\,\Omega$] (7,3.5);
  % rama inferior: R6 y eps2 (+ a la izquierda)
  \draw (0,0) -- (3.5,0) to[R, l=$R_6$, a=$12\,\Omega$] (5.4,0)
                to[battery1, l_={$\varepsilon_2 = 24\,\text{V}$}] (7,0);
  \bplus{5.95,0.32}
  \node[circ] at (3.5,3.5) {}; \node[circ] at (3.5,0) {};
\end{circuitikz}\\[2pt]
Figura 1
\end{center}
\end{minipage}

%==================================================================
\item \begin{minipage}[t]{\linewidth}
En la red de la Fig.~2, calcule la diferencia de potencial $V_P - V_Q$ entre los puntos $P$ y $Q$.
\resp{$V_P - V_Q = 8$~[V]}
\begin{center}
\begin{circuitikz}[every node/.style={font=\small}]
  % izquierda: eps2 (+ abajo, hacia P)
  \draw (0,0) to[battery1, l={$\varepsilon_2 = 14\,\text{V}$}] (0,3.5);
  \bplus{0.32,1.5}
  % superior: R1, Q, R2 y eps1 (+ a la izquierda)
  \draw (0,3.5) to[R, l=$R_1$, a=$2\,\Omega$] (3.5,3.5)
                to[R, l=$R_2$, a=$6\,\Omega$] (5.4,3.5)
                to[battery1, l={$\varepsilon_1 = 20\,\text{V}$}] (7,3.5);
  \bplus{5.95,3.18}
  % central y derecha
  \draw (3.5,0) to[R, l=$R_3$, a=$12\,\Omega$] (3.5,3.5);
  \draw (7,0) to[R, l=$R_4$, a=$5\,\Omega$] (7,3.5);
  % inferior
  \draw (0,0) -- (3.5,0) to[R, l=$R_5$, a=$1\,\Omega$] (7,0);
  \node[circ] at (3.5,0) {};
  \node[ocirc] at (3.5,3.5) {}; \node[above] at (3.5,3.6) {$Q$};
  \node[ocirc] at (0,0) {};     \node[below left] at (0,0) {$P$};
\end{circuitikz}\\[2pt]
Figura 2
\end{center}
\end{minipage}

%==================================================================
\item \begin{minipage}[t]{\linewidth}
Determine la potencia disipada en la resistencia $R_1$ del circuito de la Fig.~3.
\resp{$P_1 = 1$~[W]}
\begin{center}
\begin{circuitikz}[every node/.style={font=\small}]
  % superior: eps1 (+ a la izquierda), R1 y cable hasta la esquina
  \draw (0,3.5) to[battery1, l={$\varepsilon_1 = 3\,\text{V}$}] (1.6,3.5)
                to[R, l=$R_1$, a=$4\,\Omega$] (3.5,3.5) -- (7,3.5);
  \bplus{0.55,3.18}
  % izquierda, central y derecha (R5 abajo, R2 arriba)
  \draw (0,0) to[R, l=$R_3$, a=$8\,\Omega$] (0,3.5);
  \draw (3.5,0) to[R, l=$R_4$, a=$12\,\Omega$] (3.5,3.5);
  \draw (7,0) to[R, l=$R_5$, a=$1\,\Omega$] (7,1.75)
              to[R, l=$R_2$, a=$3\,\Omega$] (7,3.5);
  % inferior: eps2 (+ a la derecha) y R6
  \draw (0,0) -- (3.5,0);
  \draw (5.2,0) to[battery1, l={$\varepsilon_2 = 12\,\text{V}$}] (3.5,0);
  \bplus{4.6,0.32}
  \draw (5.2,0) to[R, l=$R_6$, a=$8\,\Omega$] (7,0);
  \node[circ] at (3.5,3.5) {}; \node[circ] at (3.5,0) {};
\end{circuitikz}\\[2pt]
Figura 3
\end{center}
\end{minipage}

%==================================================================
\item \begin{minipage}[t]{\linewidth}
Para el circuito de la Fig.~4, aplique las leyes de Kirchhoff para obtener la corriente $I$ que
circula por $R_4$ en el sentido indicado por la flecha. \resp{$I = 2$~[A]}
\begin{center}
\begin{circuitikz}[every node/.style={font=\small}]
  % superior y derecha
  \draw (0,3.5) to[R, l=$R_1$, a=$5\,\Omega$] (3.5,3.5)
                to[R, l=$R_2$, a=$4\,\Omega$] (7,3.5) -- (7,0);
  % izquierda: eps1 (+ abajo) y R3
  \draw (0,0) -- (0,0.3) to[battery1, l={$\varepsilon_1 = 4\,\text{V}$}] (0,1.9)
              to[R, l=$R_3$, a=$2\,\Omega$] (0,3.5);
  \bplus{0.32,0.85}
  % central: R4 con flecha I hacia arriba
  \draw (3.5,0) to[R, l={$R_4 = 3\,\Omega$}, f>_={$I$}] (3.5,3.5);
  % inferior: eps2 (+ a la derecha), R5 y R6
  \draw (1.7,0) to[battery1, l={$\varepsilon_2 = 24\,\text{V}$}] (0,0);
  \bplus{1.1,0.32}
  \draw (1.7,0) to[R, l=$R_5$, a=$1\,\Omega$] (3.5,0)
                to[R, l=$R_6$, a=$4\,\Omega$] (7,0);
  \node[circ] at (3.5,3.5) {}; \node[circ] at (3.5,0) {};
\end{circuitikz}\\[2pt]
Figura 4
\end{center}
\end{minipage}

%==================================================================
\item \begin{minipage}[t]{\linewidth}
En el circuito de la Fig.~5, calcule la potencia que entrega la fuente $\varepsilon_1$ e indique si
\'esta entrega o absorbe energ\'ia. \resp{$P_{\varepsilon_1} = 12$~[W]} (entrega energ\'ia al circuito)
\begin{center}
\begin{circuitikz}[every node/.style={font=\small}]
  % superior
  \draw (0,3.5) to[R, l=$R_1$, a=$8\,\Omega$] (3.5,3.5)
                to[R, l=$R_2$, a=$4\,\Omega$] (7,3.5);
  % izquierda: eps1 (+ arriba)
  \draw (0,3.5) to[battery1, l_={$\varepsilon_1 = 12\,\text{V}$}] (0,0);
  \bplus{0.32,2.0}
  % central: eps2 (+ abajo) y R3
  \draw (3.5,0) to[battery1, l={$\varepsilon_2 = 5\,\text{V}$}] (3.5,1.6)
                to[R, l=$R_3$, a=$12\,\Omega$] (3.5,3.5);
  \bplus{3.82,0.55}
  % derecha e inferior
  \draw (7,0) to[R, l=$R_4$, a=$10\,\Omega$] (7,3.5);
  \draw (0,0) -- (3.5,0) to[R, l=$R_5$, a=$2\,\Omega$] (7,0);
  \node[circ] at (3.5,3.5) {}; \node[circ] at (3.5,0) {};
\end{circuitikz}\\[2pt]
Figura 5
\end{center}
\end{minipage}

%==================================================================
\item \begin{minipage}[t]{\linewidth}
Para la red de la Fig.~6, determine la diferencia de potencial $V_P - V_Q$ entre los puntos $P$ y $Q$.
\resp{$V_P - V_Q = 3$~[V]}
\begin{center}
\begin{circuitikz}[every node/.style={font=\small}]
  % superior e izquierda
  \draw (0,3.5) -- (3.5,3.5) to[R, l=$R_1$, a=$5\,\Omega$] (7,3.5);
  \draw (0,0) to[R, l=$R_3$, a=$2\,\Omega$] (0,3.5);
  % central: eps1 (+ arriba)
  \draw (3.5,3.5) to[battery1, l_={$\varepsilon_1 = 2\,\text{V}$}] (3.5,0);
  \bplus{3.82,2.0}
  % derecha: R4 abajo, R2 arriba
  \draw (7,0) to[R, l=$R_4$, a=$6\,\Omega$] (7,1.75)
              to[R, l=$R_2$, a=$10\,\Omega$] (7,3.5);
  % inferior: R5 y eps2 (+ a la izquierda)
  \draw (0,0) -- (3.5,0) to[R, l=$R_5$, a=$1\,\Omega$] (5.4,0)
                to[battery1, l_={$\varepsilon_2 = 20\,\text{V}$}] (7,0);
  \bplus{5.95,0.32}
  \node[circ] at (3.5,3.5) {}; \node[circ] at (3.5,0) {};
  \node[ocirc] at (0,0) {};   \node[below left] at (0,0) {$P$};
  \node[ocirc] at (7,3.5) {}; \node[above right] at (7,3.5) {$Q$};
\end{circuitikz}\\[2pt]
Figura 6
\end{center}
\end{minipage}

%==================================================================
\item \begin{minipage}[t]{\linewidth}
Calcule la potencia total disipada en el conjunto de todas las resistencias del circuito de la
Fig.~7. \resp{$P_{\text{tot}} = 38$~[W]}
\begin{center}
\begin{circuitikz}[every node/.style={font=\small}]
  % superior y derecha
  \draw (0,3.5) to[R, l=$R_1$, a=$2\,\Omega$] (3,3.5)
                to[R, l=$R_2$, a=$10\,\Omega$] (6,3.5) -- (8.5,3.5) -- (8.5,0);
  % izquierda: eps1 (+ abajo)
  \draw (0,0) to[battery1, l={$\varepsilon_1 = 20\,\text{V}$}] (0,3.5);
  \bplus{0.32,1.5}
  % central
  \draw (6,0) to[R, l=$R_3$, a=$2\,\Omega$] (6,3.5);
  % inferior: R4, R5 y eps2 (+ a la izquierda)
  \draw (0,0) to[R, l=$R_4$, a=$6\,\Omega$] (3,0)
              to[R, l=$R_5$, a=$6\,\Omega$] (6,0)
              to[battery1, l_={$\varepsilon_2 = 8\,\text{V}$}] (8.5,0);
  \bplus{7.0,0.32}
  \node[circ] at (6,3.5) {}; \node[circ] at (6,0) {};
\end{circuitikz}\\[2pt]
Figura 7
\end{center}
\end{minipage}

%==================================================================
\item \begin{minipage}[t]{\linewidth}
En el circuito de la Fig.~8, use las leyes de Kirchhoff para determinar la corriente $I$ en $R_7$
(sentido de referencia indicado por la flecha) e interprete su signo.
\resp{$I = -1$~[A]}: la corriente circula en sentido contrario a la flecha.
\begin{center}
\begin{circuitikz}[every node/.style={font=\small}]
  % superior: eps1 (+ a la izquierda) y R1
  \draw (0,5) to[battery1, l={$\varepsilon_1 = 15\,\text{V}$}] (3.5,5)
              to[R, l=$R_1$, a=$3\,\Omega$] (7,5);
  \bplus{1.5,4.68}
  % izquierda: R4 y R2; abajo eps3 (+ arriba)
  \draw (0,1.7) to[R, l=$R_4$, a=$3\,\Omega$] (0,3.35)
                to[R, l=$R_2$, a=$3\,\Omega$] (0,5);
  \draw (0,1.7) to[battery1, l_={$\varepsilon_3 = 9\,\text{V}$}] (0,0);
  \bplus{0.32,1.1}
  % central: eps2 (+ arriba) y R5
  \draw (3.5,5) to[battery1, l_={$\varepsilon_2 = 10\,\text{V}$}] (3.5,3.5) -- (3.5,1.7);
  \bplus{3.82,4.5}
  \draw (3.5,0) to[R, l=$R_5$, a=$1\,\Omega$] (3.5,1.7);
  % derecha: R3
  \draw (7,0) -- (7,3.3) to[R, l=$R_3$, a=$3\,\Omega$] (7,5);
  % inferior: R6 y R7 (flecha I hacia la derecha)
  \draw (0,0) to[R, l=$R_6$, a=$8\,\Omega$] (3.5,0)
              to[R, l={$R_7 = 2\,\Omega$}, f>_={$I$}] (7,0);
  \node[circ] at (3.5,5) {}; \node[circ] at (3.5,0) {};
\end{circuitikz}\\[2pt]
Figura 8
\end{center}
\end{minipage}

%==================================================================
\item \begin{minipage}[t]{\linewidth}
Para la red de la Fig.~9, determine la diferencia de potencial $V_P - V_Q$ entre los puntos $P$ y $Q$.
\resp{$V_P - V_Q = -14$~[V]}
\begin{center}
\begin{circuitikz}[every node/.style={font=\small}]
  % superior: R1 hasta Q
  \draw (0,6.4) -- (3.5,6.4) to[R, l=$R_1$, a=$1\,\Omega$] (7,6.4);
  % izquierda superior: eps1 (+ abajo) y R2
  \draw (0,3.2) to[battery1, l={$\varepsilon_1 = 4\,\text{V}$}] (0,4.8)
                to[R, l=$R_2$, a=$5\,\Omega$] (0,6.4);
  \bplus{0.32,3.75}
  % rama central: R4, R5 y bajada hasta P
  \draw (0,3.2) to[R, l=$R_4$, a=$2\,\Omega$] (1.8,3.2)
                to[R, l=$R_5$, a=$2\,\Omega$] (3.6,3.2) -- (3.6,0);
  % izquierda inferior: R7 abajo, R6 arriba
  \draw (0,0) to[R, l=$R_7$, a=$5\,\Omega$] (0,1.6)
              to[R, l=$R_6$, a=$6\,\Omega$] (0,3.2);
  % inferior: R8 hasta P y cable hasta la esquina
  \draw (0,0) to[R, l=$R_8$, a=$1\,\Omega$] (3.6,0) -- (7,0);
  % derecha: R3 arriba y eps2 (+ arriba) abajo
  \draw (7,3.2) to[R, l=$R_3$, a=$3\,\Omega$] (7,6.4);
  \draw (7,3.2) to[battery1, l_={$\varepsilon_2 = 20\,\text{V}$}] (7,0);
  \bplus{7.32,1.85}
  \node[circ] at (0,3.2) {};
  \node[ocirc] at (3.6,0) {}; \node[below] at (3.6,-0.08) {$P$};
  \node[ocirc] at (7,6.4) {}; \node[above right] at (7,6.4) {$Q$};
\end{circuitikz}\\[2pt]
Figura 9
\end{center}
\end{minipage}

%==================================================================
\item \begin{minipage}[t]{\linewidth}
Calcule la potencia total disipada en el conjunto de todas las resistencias del circuito de la
Fig.~10. \resp{$P_{\text{tot}} = 43$~[W]}
\begin{center}
\begin{circuitikz}[every node/.style={font=\small}]
  % superior: R1, R2 y R3
  \draw (0,6.2) to[R, l=$R_1$, a=$2\,\Omega$] (2,6.2)
                to[R, l=$R_2$, a=$5\,\Omega$] (4,6.2)
                to[R, l=$R_3$, a=$8\,\Omega$] (6.5,6.2);
  % izquierda superior: R6 abajo, R4 arriba
  \draw (0,3.1) to[R, l=$R_6$, a=$3\,\Omega$] (0,4.65)
                to[R, l=$R_4$, a=$4\,\Omega$] (0,6.2);
  % izquierda inferior: R7 y eps3 (+ arriba)
  \draw (0,1.55) to[R, l=$R_7$, a=$12\,\Omega$] (0,3.1);
  \draw (0,1.55) to[battery1, l_={$\varepsilon_3 = 2\,\text{V}$}] (0,0);
  \bplus{0.32,1.0}
  % rama interior: eps1 (+ abajo)
  \draw (0,3.1) -- (3.5,3.1);
  \draw (3.5,0) to[battery1, l={$\varepsilon_1 = 18\,\text{V}$}] (3.5,3.1);
  \bplus{3.82,1.3}
  % derecha: R5 arriba y eps2 (+ arriba) abajo
  \draw (6.5,3.1) to[R, l=$R_5$, a=$10\,\Omega$] (6.5,6.2);
  \draw (6.5,3.1) to[battery1, l_={$\varepsilon_2 = 6\,\text{V}$}] (6.5,0);
  \bplus{6.82,1.8}
  % inferior
  \draw (0,0) to[R, l=$R_8$, a=$4\,\Omega$] (3.5,0) -- (6.5,0);
  \node[circ] at (0,3.1) {}; \node[circ] at (3.5,0) {};
\end{circuitikz}\\[2pt]
Figura 10
\end{center}
\end{minipage}

%==================================================================
\item \begin{minipage}[t]{\linewidth}
En el circuito de la Fig.~11, determine la potencia disipada en la resistencia $R_4$.
\resp{$P_4 = 18$~[W]}
\begin{center}
\begin{circuitikz}[every node/.style={font=\small}]
  % lazo exterior superior y bajada derecha hasta el nodo
  \draw (0,6) -- (6.5,6) -- (6.5,3);
  % izquierda: eps1 (+ abajo) arriba y eps2 (+ arriba) abajo
  \draw (0,3) to[battery1, l={$\varepsilon_1 = 24\,\text{V}$}] (0,6);
  \bplus{0.32,4.25}
  \draw (0,3) to[battery1, l_={$\varepsilon_2 = 12\,\text{V}$}] (0,0);
  \bplus{0.32,1.75}
  % derecha: eps3 (+ arriba)
  \draw (6.5,3) to[battery1, l_={$\varepsilon_3 = 4\,\text{V}$}] (6.5,0);
  \bplus{6.82,1.75}
  % rama interior: R1, R2 y R3
  \draw (2.8,3) to[R, l=$R_1$, a=$2\,\Omega$] (4.65,3)
                to[R, l=$R_2$, a=$2\,\Omega$] (6.5,3);
  \draw (2.8,0) to[R, l=$R_3$, a=$8\,\Omega$] (2.8,3);
  % inferior
  \draw (0,0) to[R, l=$R_4$, a=$2\,\Omega$] (2.8,0)
              to[R, l=$R_5$, a=$4\,\Omega$] (6.5,0);
  \node[circ] at (6.5,3) {}; \node[circ] at (2.8,0) {};
\end{circuitikz}\\[2pt]
Figura 11
\end{center}
\end{minipage}

\end{enumerate}

\end{document}
